% =============================================================
%  OpenTuring Project Cover --- public template (for \input reuse)
%  Depends: fontspec, xeCJK, xcolor, tikz, fontawesome5
%           \usetikzlibrary{positioning,calc}
% =============================================================

\definecolor{coverprimary}{HTML}{5B2D8E}
\definecolor{coveraccent}{HTML}{B03A2E}
\definecolor{coveramber}{HTML}{D9A441}
\definecolor{coverteal}{HTML}{1E8E8E}
\definecolor{coverdark}{HTML}{241A33}
\definecolor{covermuted}{HTML}{6B6478}
\definecolor{goldpanel}{RGB}{255,248,230}
\definecolor{badgealgo}{HTML}{2E5A9E}
\definecolor{badgetype}{HTML}{D9A441}
\definecolor{badgelang}{HTML}{1E8E8E}
\definecolor{badgetheory}{HTML}{C0395B}

\newcommand{\openturingslide}{%
\begin{frame}[plain]
\begin{tikzpicture}[remember picture, overlay]
  \fill[coverprimary!4]   (current page.north west) rectangle (current page.south east);
  \fill[coverdark, opacity=0.06] (current page.south west) rectangle ([yshift=0.4cm]current page.south east);
  \draw[coverprimary, opacity=0.30, line width=0.8pt] ([yshift=0.4cm]current page.south west) -- ([yshift=0.4cm]current page.south east);

  \node[anchor=center, font=\fontsize{24}{30}\selectfont\bfseries, text=coverprimary!90!black]
    at ([yshift=3.9cm]current page.center) {OpenTuring 已开放源码};
  \node[anchor=center, font=\fontsize{14}{18}\selectfont, text=coverprimary!60!black]
    at ([yshift=3.05cm]current page.center) {为历届图灵奖得主立传};

  \draw[coveraccent!45, line width=1.0pt] ([yshift=2.45cm, xshift=-5.0cm]current page.center)
    -- ([yshift=2.45cm, xshift=5.0cm]current page.center);

  \node[anchor=center, font=\fontsize{8.2}{11.5}\selectfont, text=coverdark!72, align=center, text width=13cm]
    at ([yshift=1.55cm]current page.center)
    {从 Alan Perlis、高德纳、Dijkstra，到 Hinton、Bengio、LeCun，\\[2pt]
     OpenTuring 正在为计算机科学家建立开放、持续演进的图灵奖人物档案。};

  \node[anchor=center, font=\fontsize{8}{11}\selectfont\bfseries, text=coverprimary!82!black, align=center]
    at ([yshift=0.45cm]current page.center)
    {这不只是获奖名单，而是一部由社区共同构建的现代计算机科学人物史：};

  \node[anchor=north west, font=\fontsize{7}{9.2}\selectfont, text=coverdark!78, align=left]
    at ([xshift=-6.3cm, yshift=-0.05cm]current page.center)
    {\textcolor{badgealgo}{\faBook}\enspace\textbf{历届得主} --- 持续收录、整理与补充\\[4pt]
     \textcolor{coveraccent}{\faMap}\enspace\textbf{人物生平} --- 教育经历、学术生涯、重要节点\\[4pt]
     \textcolor{badgelang}{\faTh}\enspace\textbf{核心贡献} --- 算法、语言、系统与理论};
  \node[anchor=north west, font=\fontsize{7}{9.2}\selectfont, text=coverdark!78, align=left]
    at ([xshift=0.3cm, yshift=-0.05cm]current page.center)
    {\textcolor{badgetheory}{\faTrophy}\enspace\textbf{荣誉与传承} --- 图灵奖、哥德尔奖、von Neumann 奖章等\\[4pt]
     \textcolor{coverteal}{\faGlobe}\enspace\textbf{学术谱系} --- 师承、门生、合作与思想传承\\[4pt]
     \textcolor{coveramber}{\faCode}\enspace\textbf{开放共建} --- 欢迎任何人补充、纠错、完善};

  \node[draw=coveraccent!50, fill=goldpanel, rounded corners=5pt, inner sep=9pt, text width=12.0cm, align=center]
    at ([yshift=-2.35cm]current page.center)
    {{\fontsize{8.6}{11}\selectfont\bfseries\color{coverdark!88} 计算机科学史，不应该只存在于少数人的书架上。}\\[2pt]
     {\fontsize{7.5}{9.5}\selectfont\color{covermuted} 让每一位图灵奖得主的故事都可以被发现、被理解、被连接。}};

  \node[anchor=south, font=\fontsize{7.8}{9.8}\selectfont\bfseries, text=coveraccent!80!black]
    at ([yshift=0.64cm]current page.south) {\faIcon{star}\enspace Star OpenTuring · 加入共建\enspace|\enspace https://github.com/OpenMathAI/OpenMath};
  \node[anchor=south, font=\fontsize{6.2}{7.8}\selectfont, text=covermuted]
    at ([yshift=0.38cm]current page.south) {欢迎 Star、Issue、PR，共同建设一部开放的图灵奖人物史。};
\end{tikzpicture}
\end{frame}
}
